In passive potential-flow networks, such as gas and water networks
without active elements, each nomination and choice of edge-law
coefficients determines a unique physical state. Optimizing a pressure
difference or arc flow of this state over uncertain nominations and
coefficients is a nonconvex problem arising in robust validation and
design. We determine how its complexity in the Turing bit model depends
on topology, observable, uncertainty model, and required output. At
fixed maximum cycle rank per biconnected block, the largest prescribed
potential difference over a balanced nomination box and independent
resistance intervals is computable to additive accuracy $\eps$ in time
polynomial in the input length and $\log(1/\eps)$, with an exactly
admissible rational scenario; arc-flow extrema are computable exactly,
so robust capacity validation is exact. On cacti, this answers the
additive version of an open question recorded in a recent survey, while exact
pressure comparison contains Square-Root Sum and, for one source and
sink, is equivalent to it. For quadratic laws and fixed nominations,
interval hulls of independent resistance sets preserve all extrema
precisely on trees for weighted potentials, cacti for potential
differences and linear flows, and series-parallel graphs for arc flows.
Over nomination boxes, weighted
potentials with growing support are NP-hard on trees, yet admit
accuracy-bit algorithms in three structured settings with fixed
nomination dimension or support. On cacti with fixed nominations, exact
capacity recovery survives global coefficient correlations, which
nevertheless make total flow and minimum dissipation strongly NP-hard;
maximum dissipation over a resistance polytope is tractable on every
connected graph. Rational energy, Bregman, and curvature certificates
turn numerical solutions into exactly checkable guarantees.
